% =====================================================================
% sections/discussion.tex — §5.10 综合讨论与本章小结（rev.1，2026-07-07 重流程落地）
% 依据：已批方案 paper/thesis_ch5/section_5_10_discussion_writing_plan.md
%   （2026-07-07 用户"同意推荐"，§0 六项裁决全按推荐锁定）。落地按方案 §2 逐子节
%   论点骨架 + §4 claim→锚点映射表（C0.1–C5.4 逐条核锚）+ §5 红线自查表执行。
% 六项已锁裁决：① 开篇段（不编号）+ 5 正文子节，无 §5.10.m（方案 A；spec 指定给
%   .m 的三资产由正文承载——β1 四轴表进 §5.10.2、特权四态表进 §5.10.3、五重限定
%   与 limitations 散文进 §5.10.4）；② 节标题「综合讨论与本章小结」、label
%   sec:ch5_discussion；③ 两表零图、两表零新数字（定性差异 + 节内指针；唯一携带的
%   定量签名"约 6 个百分点、约合差距之半"系 §5.6.4 章内已有数字并附锚）；
%   ④ sim2real 与 limitations 合并为 §5.10.4；⑤ 生物回扣一句置 §5.10.5 末段
%   （四护栏：仅一句 / 贡献句无生物专名 / "功能上的替代，而非结构上的仿制"字样 /
%   不新增 claim）；⑥ 贡献类型正面声明置 §5.10.5 第二段（三类型；正向声明、
%   不 itemize、禁反向宣告）。
% 排版落地说明：表 5.10-1 按"维度 × 两协议"三列呈现（方案 §3 定案的"两行 × 多列"
%   之转置；内容逐格一致，转置属排版层落地自由度，六列宽表在 138mm 版心内不可读）。
% 收束引用逐字对齐（起草时开原文对照，均已核）：
%   伏笔句 §5.8.5"特权信息在两条线索、多种工况下呈现的不同失效方式，其统一收束
%   留待本章末节"/ §5.9.2"两个最优取值在同一机制下的统一解释留待本章末节"/
%   §5.9.5"算法与数据条件交互的统一收束连同行为约束强度的跨线解释一并留待本章
%   末节"；统一句逐字对齐 intro ¶7"都在固定部署接口下提高既有轨迹中可被策略提取
%   的信息量"；中心命题原话回归 = spec §0.1 校准锚（小结 ¶1）；"跨种子稳健性优选
%   而非单点最高值优选"（§5.7.5）；"训练期特权观测的必要性随算法对既有信息的利用
%   程度而消长"（§5.7.5）；"只宜读作机制方向上的提示"（§5.6.4）；"在线情形下靠
%   延长时序窗闭合差距的途径，未必能直接迁移到离线"（§5.5.5）；§5.9.5 rev.2 质量
%   区间口径（中段去噪参照占优 / 饱和高位直接模仿至少相当）；"确有其效，但仅在
%   数据质量落于其去噪能力有事可做的区间内"（§5.9.5）；两读法并存 + 两种子限定
%   连同引用（§5.8.5）。
% 红线遵守清单（方案 §5；grep 自查通过）：零新数字（正文与两表无章内未有数字，
%   章内已有数字除"约 6 个百分点"定量签名外一律 \S\ref 指回）；§0.4 五条——LN 仅
%   一句正交定位不进清单、+23–32pp 零出现、三项失效方式各异（§5.10.3 结构执行）、
%   LN n=2 不排序、交互只方向不跨源幅度比（两表表注 + §5.10.2/§5.10.3 正文各声明
%   一次）；不写"某线调错"；禁 teacher / 裸写 ReBRAC / 「」角引号 / 直引号 / 其一
%   其二 / 信息论 / "不可能"字样 / regime / 导航 / seed 44 具名 / 报告体 marker；
%   小结零数字零新文献不开新论点；生物回扣仅一句、贡献句无生物专名；零跨章；
%   本节无新增 cite。
%
% rev.2（2026-07-08）：独立对抗复审落地（重流程第四步；findings 存档
%   section_5_10_review_findings.md；方案 §4 映射表 24 条逐条对账 + 九节接口
%   原文与 ground truth 全量实时回查，未采信 rev.1 头注自述）。两表零漂移、
%   收束引用零走样、伏笔兑现零扩权、红线 grep 零踩线；修订实质项
%   （0 CRIT / 3 HIGH / 7 MED / LOW 顺手 2 条）：
%   [HIGH-scope] §5.10.4 "与本章全部正面结论一样，限定于五重设定之内"假全称
%     （§5.5 时序闭合 / §5.8.1 锚点 / §5.9 判别证据均在到达导向甚至临界工况）
%     → 收窄为"与离线主线的全部正面结论一样"（§5.7.5 自划边界同此）。
%   [HIGH-综合句] §5.10.3 四态统一句"无论在线或离线、亚临界或临界"——策略侧
%     瓶颈分支仅态一/态四两处临界证据，无亚临界实例 → 删"、亚临界或临界"。
%   [HIGH-扩权] §5.10.4 "离线数据采集与训练可以完全建立在部署可得的传感接口
%     上"——追平单元数据由信息口径强于单点的采集器采集（§5.6.3 明言）、临界
%     数据由读全场采集器采集；paper 1 §6.5 素材只 claim 训练管线 → 改"策略的
%     训练与部署由此完全建立在部署可得的传感接口上"。
%   [MED×7] §5.10.1 机制句 hedge 回归（"便能"→"便可能"，对齐 online rev.2 假说
%     措辞）；"高品质确定性示范"→"高质量确定性采集数据"+ 表格"确定性示范"→
%     "确定性采集"（"示范"章内专属 §5.8/§5.9 全场走廊参照数据；"无注入噪声"经
%     采集脚本默认值独立核实成立）；β1 最优点决定因素删"数据质量所处区间"
%     （质量区间决定算法排序而非锚定最优点，§4.0.10 明记 β1=1≥4 跨质量档一致）；
%     "张力的第一层消解"孤悬序数→"张力首先消解于协议差异"；两读法之一补
%     "$k=4$ 观测构成下的"限定（§5.8.5 原文）；表 5.10-1 种子组格"固定五种子组"
%     →"固定种子组（筛查三种子、正式五种子）"（稳健性优选发生在三种子筛查网格，
%     与证据锚对齐）；"第一条/第二条途径、第二途径"序数指称三处溶解为具名指称。
%   [LOW 顺手] "该协议与采集器差距之半"先行词→"可部署协议"（正文+表格两处）；
%     "以横流穿越为主"→"与横流穿越"（grep 全章任务几何均为横流穿越，"为主"与
%     "各取其一"相抵）。余 LOW×3（"仿真器方便提供"语感 / 表 5.10-2 态一行名 /
%     机体↔船体章级术语分歧——本节从 setup.tex 定义之家，跨节统一留收尾统稿轮）
%     记录 findings §4。门控：CRIT+HIGH=3 < §5.9 基线 4，判定与呈报见 findings §6。
% =====================================================================

\section{综合讨论与本章小结}\label{sec:ch5_discussion}

全章的证据至此已经陈列完毕，各结果节也已就近界定了自身结论的含义与边界；本节不引入新的实验，所论全部依据前九节已呈证据。收束沿中心命题的正反两轴展开：正向轴把在线情形的时序利用与离线情形的数据组织归并为同一种信息提取，反向轴把三种为系统增添能力的手段按各自的失效方式收拢。前文留待本章末节的问题也在此逐一兑现——行为约束强度的两个最优取值在同一机制下的统一解释、算法与数据条件交互的统一收束，以及特权信息在两条线索、多种工况下不同失效方式的统一图景——本章小结最终回到 \S\ref{sec:ch5_intro} 提出的中心命题。

\subsection{正向轴的统一机制：时序利用与数据组织同属信息提取}\label{subsec:ch5_discussion_positive}

正向轴的时序利用途径在在线情形获得了受控证据。临界工况下，可部署单点配置的成功率随观测历史长度单调上升，闭合到空间参照配置的水平并贴近评估集的经验上界（\S\ref{subsec:ch5_online_ablation}、\S\ref{subsec:ch5_online_robustness}）；其机制解释是，单点流速随时间的变化编码了主导涡街脉动的相位，时序窗覆盖脱涡周期的相当比例后，策略便可能从时序节拍中恢复原本似乎须由空间探针提供的相位信息。这一途径不增添任何部署期不可得的传感，却把可部署性能推到该评估集对单点观测所允许的上界附近。

该途径的离线边界须随其结论一并陈述。本章离线各线的数据均以 $k=4$ 观测构成采集，离线学习只能复用既采数据的观测构成，\S\ref{subsec:ch5_online_implications} 因而指出，在线情形下靠延长时序窗闭合差距的途径，未必能直接迁移到离线；\S\ref{subsec:ch5_boundary_ceiling} 所划的临界工况离线上限，也正是 $k=4$ 观测构成下的经验上限。时序利用在离线一侧的可行性因此是正向轴留下的开放前沿，而非已经闭环的结论。

正向轴的另一条途径——使模仿目标与数据质量相适配——由离线两族数据的判别证据支撑。这一途径的前半在离线主线得到正面例证：高质量确定性采集数据上，以原始数据动作为模仿目标工作良好，双侧行为约束把这份数据利用到与特权协议参照相当的水平（\S\ref{subsec:ch5_rebrac_implications}）。后半由算法比较给出判别形态：数据质量处于尚未逼近上限的中段区间时，流匹配去噪的参照动作在同源对照中占优；质量高、性能逼近数据上限时，原始动作本身就是好目标，恰当加权的直接模仿与之至少相当（\S\ref{subsec:ch5_algo_implications}）。最优的模仿目标随数据条件移动，适配本身、而非任何单一固定目标，构成这条途径的内容。

离线基线暴露的数据支持结构问题与双侧约束对它的消化（\S\ref{subsec:ch5_td3bc_support}、\S\ref{subsec:ch5_rebrac_scale}），按 \S\ref{sec:ch5_intro} 的既定口径同属离线数据的有效组织：决定离线性能的不是样本计数，而是数据被组织进支持集的方式与算法对该支持的利用程度。两条途径由此在机制层面汇合——正如引言所预设的，二者实质上都在固定部署接口下提高既有轨迹中可被策略提取的信息量：时序利用在时间维度上展开单点信号，目标适配与数据组织则在离线数据内部重新划分可被可靠模仿与自举的成分。两条途径的证据形态并不对称：时序利用有在线受控消融的正面证据、其离线可行性尚属开放，目标适配有离线两族数据的判别证据、但未在在线情形检验。本小节的统一是机制口径的统一，不抹平这一证据不对称。

\subsection{行为约束强度的跨线统一解释}\label{subsec:ch5_discussion_anchor}

两条离线线索在行为约束强度上留下一处表面张力。离线主线在其协议下锁定 $\beta_1=4.0$，并以之取得全部主结果（\S\ref{subsec:ch5_rebrac_implications}）；算法比较的消融却显示，$\beta_1=1.0$ 在干净与含噪两轴同时更好（\S\ref{subsec:ch5_algo_anchor}）。若不加解释，这两个最优取值会被读作实验之间的矛盾；本小节给出 \S\ref{subsec:ch5_algo_anchor} 留下的统一解释。

张力首先消解于协议差异。两套协议在奖励设定、采集器、动作噪声与随机种子组四个维度上系统不同，表~\ref{tab:ch5_discussion_beta_axes} 逐轴并置两侧的取值与章内证据锚；其中任何一轴都可能移动锚定强度的最优点，四轴同时改变则使两个取值之间不存在可直接相减的对照。

\begin{table}[t]
\centering
\caption{行为约束强度两处最优取值所处协议的逐轴对照。离线主线与算法比较两套协议在四个维度上系统不同，其中任何一轴都可能移动锚定强度的最优点；两侧证据各出自其自身协议，本表只支撑机制方向的读数，不支撑两侧效应幅度的比较，定量细节以章内证据锚为准。}
\label{tab:ch5_discussion_beta_axes}
\small
\setlength{\tabcolsep}{5pt}
\begin{tabular}{l p{47mm} p{52mm}}
\toprule
\textbf{维度} & \textbf{离线主线（\S\ref{sec:ch5_rebrac}）} & \textbf{算法比较（\S\ref{sec:ch5_algo_compare}）} \\
\midrule
奖励设定 & 历史效率导向 & 到达导向 \\
采集器 & 横流补偿参照（确定性规则控制器） & 全场走廊参照（读取全场流场） \\
动作噪声 & 确定性采集，无注入噪声 & 无注入与 $\sigma=0.5$ 注入两档对照 \\
随机种子组 & 固定种子组（筛查三种子、正式五种子），含难例种子 & 两种子组，不含难例种子 \\
最优取值与选择准则 & $\beta_1=4.0$；跨种子稳健性优选 & $\beta_1=1.0$；干净与含噪两轴对照下的优选 \\
章内证据 & 表~\ref{tab:ch5_rebrac_screen}（按筛查协议解读） & 图~\ref{fig:ch5_fql_noise_axis} \\
\bottomrule
\end{tabular}
\end{table}

种子组一轴的后果最直接，因为它改变的不只是数据，还有选择准则本身。主线取值是跨种子稳健性优选而非单点最高值优选：筛查网格中锚定较弱的配置不乏更高的单种子成绩，但难例种子在弱锚定下系统性偏弱，这一证据由表~\ref{tab:ch5_rebrac_screen} 按筛查协议在章内承载（\S\ref{subsec:ch5_rebrac_implications}）。算法比较一侧的种子组不含这一难例种子，为稳健性抬高锚定的动机随之不在，含噪数据上弱锚定的优势便直接显现（\S\ref{subsec:ch5_algo_anchor}）。

由此得到的统一解释把表面矛盾化为一项发现：锚定强度在采集器、奖励与动作噪声共同张成的条件空间里没有单一全局最优，其最优点随模仿目标是否含噪与协议的稳健性要求移动，两个取值各是其自身协议下的合理优选。这正是模仿目标与数据质量相适配这一正向途径在同一算法内部的再现——模仿目标与数据条件的适配决定超参数的最优点，一如它决定算法之间的相对排序；\S\ref{subsec:ch5_algo_anchor} 的机制表述“其取值由模仿目标是否含噪决定，与算法家族无关”，与 \S\ref{subsec:ch5_algo_crosssource} 中该机制随数据源保持的证据，共同支撑这一读法。与表~\ref{tab:ch5_discussion_beta_axes} 表注一致，这一跨线对照只作机制方向的读数：两侧效应出自不同协议，其幅度不作比较。

\subsection{反向轴的收束：三种增添能力手段的失效方式各异}\label{subsec:ch5_discussion_negative}

反向轴的三种手段——更多空间传感器、向价值网络注入特权观测、采用更具表达力的策略先验——在本章证据中各有其精确的失效形式。把它们收束到中心命题之下，恰恰要求保持这些失效方式的区别，而非把它们混为一律无用。

空间传感器一项的界定最为清晰：增设前向探针确实闭合临界工况下的差距，但该探针在部署中不可得，且其提供的相位信息可由对单点观测的充分时序利用替代（\S\ref{subsec:ch5_online_implications}）。它的失效形式因此不是无效，而是非必需——有效、不可部署、可替代，这三个限定共同构成其收束表述。

向价值网络注入特权观测在全章出现于四种互不相同的协议之下，表~\ref{tab:ch5_discussion_priv_states} 把四处证据的作用形态并置。在线临界工况下，它不闭合可部署配置与空间参照之间的差距，收敛峰值跨重复训练稳定锁定在低位（\S\ref{subsec:ch5_online_ablation}）。离线亚临界工况的单侧约束基线上，它呈现方向性的有益：点估计约 $6$ 个百分点、约合可部署协议与采集器差距之半，但显著性未确立，只宜读作机制方向上的提示（\S\ref{subsec:ch5_td3bc_critic}）。同一数据上的双侧约束主线不再出现可检测的均值抬升，特权观测的残余价值收缩为对单一难例种子的救援（\S\ref{subsec:ch5_rebrac_worldcomp}）。而离线临界工况下，它连挽救完全失效的能力也没有，价值侧成因被排除——引用这一临界结论须连同其两种子限定（\S\ref{subsec:ch5_boundary_asym}、\S\ref{subsec:ch5_boundary_implications}），且零成功同时兼容 $k=4$ 观测构成下的表征不足与非对称机制未把特权信息转化给策略两种读法，本章不在两读法之间裁决。

\begin{table}[t]
\centering
\caption{训练期向价值网络注入特权观测在全章四种协议下的作用形态收束。四行出自线别、工况、算法与奖励设定各异的四套协议，本表收束的是作用形态的方向图景，不构成受控比较，各行之间的效应幅度不可互相并排；全部定量细节以章内证据锚为准。}
\label{tab:ch5_discussion_priv_states}
\small
\setlength{\tabcolsep}{4pt}
\begin{tabular}{p{34mm} p{44mm} p{36mm} l}
\toprule
\textbf{工况与协议} & \textbf{作用形态} & \textbf{证据强度限定} & \textbf{章内证据} \\
\midrule
在线临界；标准 SAC，到达导向奖励 & 不闭合差距；收敛峰值跨重复训练稳定锁定在低位 & 两随机种子配对复现 & \S\ref{subsec:ch5_online_ablation} \\
\addlinespace
离线亚临界；单侧约束基线 TD3+BC，历史效率导向奖励 & 方向性有益：点估计约 $6$ 个百分点、约合可部署协议与采集器差距之半 & 五随机种子；显著性未确立，只宜读作机制方向提示 & \S\ref{subsec:ch5_td3bc_critic} \\
\addlinespace
离线亚临界；双侧约束主线 ReBRAC-Q，历史效率导向奖励 & 均值无可检测抬升；残余为对单一难例种子的救援 & 五随机种子 & \S\ref{subsec:ch5_rebrac_worldcomp} \\
\addlinespace
离线临界；ReBRAC-Q，到达导向奖励，高品质示范数据 & 不能挽救完全失效；价值侧成因被排除 & 两随机种子；零成功为零方差强信号，引用须连同种子限定 & \S\ref{subsec:ch5_boundary_asym} \\
\bottomrule
\end{tabular}
\end{table}

四种作用形态贯穿着同一个条件化结构：训练期特权观测的必要性随算法价值侧的自足程度与工况难度而变。价值网络自身的改进信号不足时，补入特权观测呈方向性有益；价值侧的稳定化已由不依赖特权信息的支持惩罚完成时，它退为非必需——\S\ref{subsec:ch5_rebrac_implications} 所述“训练期特权观测的必要性随算法对既有信息的利用程度而消长”正是这一结构的主线表述；而当瓶颈落在策略侧可获取的信息时，无论在线或离线，向价值网络增添信息都不再改变可部署性能。四态出自线别、工况、算法与奖励设定各异的四套协议，这一收束因此是作用形态方向图景的收束，不构成受控比较，各态之间的幅度不可互相并排（表~\ref{tab:ch5_discussion_priv_states} 表注同此声明）。

更具表达力的策略先验一项以不普遍有效的精确形式收束：生成式先验确有其效，但仅在数据质量落于其去噪能力有事可做的区间内（\S\ref{subsec:ch5_algo_implications}）；质量饱和的高位上，它不再提供恰当加权的直接模仿所没有的东西。与这三项相区别，价值网络的层归一化是表征稳定性的必要基础设施，与“增添能力是否必需”这一命题维度正交，不属于本节收束的清单（\S\ref{subsec:ch5_rebrac_ln}）。三种手段各自的失效形式——有效但非必需、条件化的非必需以至无济于事、不普遍有效——合起来正是中心命题反向轴的原话：或非必需，或不普遍有效。

\subsection{对部署的含义与本章结论的适用范围}\label{subsec:ch5_discussion_deployment}

本章对部署最直接的含义来自离线主线：在本章考察范围内，可部署离线训练不依赖任何部署期不可得的信号，即可达到与训练期接收特权观测的同类协议在统计意义上持平的水平（\S\ref{subsec:ch5_rebrac_worldcomp}）。对真实平台而言，这意味着离线训练管线无需仿真器方便提供的特权通道：部署期不可得的等效来流估计，既不必在训练期提供给价值网络，也不出现在部署回路中，策略的训练与部署由此完全建立在部署可得的传感接口上（\S\ref{subsec:ch5_method_protocol}）。这一含义断言的是统计意义上的持平而非超越，其成立范围由下文的限定共同划定。

上述含义与离线主线的全部正面结论一样，限定于五重设定之内：$s_0$ 单点传感接口、圆柱卡门涡街尾流场、REMUS-100 平台、亚临界流动工况，与历史效率导向奖励。其中奖励一重已有锚点扩展：亚临界锚点单元显示，主线的可部署结论在到达导向奖励下保持，但该对照跨越奖励设定、种子数与评估回合数多个口径，只作量级参考（\S\ref{subsec:ch5_boundary_anchor}）。数据规模退化的翻转（\S\ref{subsec:ch5_rebrac_scale}）与上述部署含义同受这组边界约束，本章不把它们外推为一般性的算法结论。

本章结论的适用范围还受几类章级限制约束。全部结论出自仿真，未经实航验证；平台、尾流场族与任务几何各取其一——REMUS-100、圆柱卡门涡街与横流穿越——跨平台与跨流场族的稳健性不在本章证据范围内。评估与数据采集取自同一任务分布，分布偏移下的稳健性未测；离线数据规模位于千回合量级，数据极充足时双侧行为约束是否仍然必要，本章未验。特权观测仅为二维的沿机体采样等效来流，\S\ref{subsec:ch5_discussion_negative} 的四态收束因此限于这一最窄的特权设置，更高维特权信号下的图景本章未回答。

统计形态的限制同样须在收束层保持。全章种子规模沿五、三、二的阶梯递减，临界边界与个别消融是两种子的存在性证据，其证据分量已在 \S\ref{subsec:ch5_boundary_implications} 与 \S\ref{subsec:ch5_algo_implications} 各自写明；本节的收束不解除任何一处限定，引用相关结论时应连同其种子规模。离线各线固定 $k=4$ 观测构成，时序利用途径在离线一侧的可行性因此悬而未决——它既是 \S\ref{subsec:ch5_discussion_positive} 所述证据不对称的来源，也是本章正向轴留给后续工作的最大开放前沿。

\subsection{本章小结}\label{subsec:ch5_discussion_summary}

本章围绕一个中心命题展开检验：在部署约束下，提升性能的关键在于用好已有的信息与数据，而非为系统增添能力。前者落实为两条途径——在时间维度上充分利用单点观测，以及使模仿目标与数据质量相适配；后者——更多空间传感器、向价值网络注入特权观测、采用更具表达力的策略先验——则或非必需，或不普遍有效。前九节的证据支持这一命题在本章设定内成立：时序利用经在线临界工况的受控消融确证，模仿目标的质量适配在离线两族数据上获得判别证据，数据组织的作用由离线基线与主线的对照界定；三种增添能力的手段各以其精确的失效形式划界。命题的适用范围由亚临界与临界工况之间的边界，连同上一小节的各重限定共同给出。

本章的贡献可归为部署约束下信息利用的一组经验规律与设计准则：单点观测的时序访问决定在线可部署性能能否闭合差距，模仿目标与数据质量的适配决定离线算法及其锚定强度的相对表现，训练期特权观测的必要性随算法与工况条件化。支撑这些规律的受控判别框架——单变量消融、协议对照与预登记判读——连同其中经受检验的 Q 归一化双侧行为约束方法，构成可为后续工作复用的方法制品。

本章对空间与时间两个维度的判别，也为引言中的生物图景给出一个工程侧的回应：自然界以沿体表分布的侧线阵列在空间上弥补无法测量全场，单点传感的航行器则可在时间维度上重组同一相位信息——这是功能上的替代，而非结构上的仿制。在博士论文的整体论证中，本章据此提供了部署约束下单点感知运动规划的可行性证明、失效边界与算法决定因素；其用好已有信息与数据的准则不依赖任何特定算法，可延伸至同类部署受限的自主系统。
